% Auto-generated by scripts/latex_synthesis.py -- do not hand-edit; regenerate from the
% source synthesis/<slug>.md instead, since edits here will be overwritten on the next run.
%
% Thesis-chapter style: report class (numeric \cite by default, no natbib/biblatex needed),
% standard margins, no title page/TOC -- this is meant to be read as a single chapter, not a
% standalone thesis. See scripts/latex_synthesis.py's CITATION REPLACEMENT STRATEGY comment
% for how the "Author et al. (Year)" text was turned into \cite{} calls below.
\documentclass[12pt]{report}
\usepackage[margin=1in]{geometry}
\usepackage[utf8]{inputenc}
\usepackage[T1]{fontenc}

\begin{document}

\chapter{how do these papers handle attention/alignment mechanisms}

% Cross-paper synthesis generated by scripts/synthesize.py.

Both papers address the problem of relating positions across a sequence, but they differ sharply in how central attention is to the overall architecture. \cite{1409.0473} introduce attention as a targeted fix to a specific bottleneck: the fixed-length context vector of the RNN encoder-decoder degrades on long sentences, so they add a learned alignment model that lets the decoder compute a weighted combination of bidirectional-RNN annotations at each output step, effectively replacing a single fixed vector with a per-timestep soft-searched context \cite{1409.0473}. Attention here is a supplementary mechanism layered on top of recurrence — the encoder and decoder are still RNNs, and the alignment weights are produced by a small feedforward network scoring decoder state against each source annotation \cite{1409.0473}. \cite{1706.03762} take the opposite stance, arguing that if attention already lets a model draw dependencies between arbitrary positions regardless of distance, then recurrence is not merely optional but actively limiting, since it prevents within-sequence parallelization and caps computation to a fixed number of operations only in the attention case \cite{1706.03762}. The Transformer accordingly discards recurrence entirely and relies on scaled dot-product self-attention, computed in parallel across multiple learned projections ("multi-head attention") for encoder self-attention, decoder self-attention, and encoder-decoder attention alike \cite{1706.03762}.

This contrast reflects a broader trajectory from attention-as-addition to attention-as-foundation. Bahdanau et al.'s alignment mechanism demonstrates empirically that soft, differentiable alignment outperforms compressing a whole sentence into one vector, and that the resulting attention weights correspond to linguistically plausible word alignments \cite{1409.0473}. Vaswani et al. generalize this insight into a full alignment substitute for sequential computation, showing that a purely attention-based architecture not only matches but exceeds RNN-based translation quality while training faster and generalizing to non-translation tasks such as constituency parsing \cite{1706.03762}. Where Bahdanau et al.'s alignment model is a single additive-attention computation between decoder state and encoder annotations, Vaswani et al.'s scaled dot-product formulation is explicitly framed as a faster, more space-efficient alternative to additive attention, with multi-head projection compensating for the loss of resolution that a single attention function would otherwise incur \cite{1706.03762}.

\bibliographystyle{plain}
\bibliography{how-do-these-papers-handle-attention-alignment-mechanisms}

\end{document}
